\chapter{EDA III: Outliers, Missing Values and Leaky Variables, in Depth}
\companion{StatQuest: The Normal Distribution, Clearly Explained (search)}{\yts{StatQuest+The+Normal+Distribution+Clearly+Explained}}

Chapter 13 introduced missing data, outliers and imbalance from a modelling angle. The checklist asks for more detail: the 68--95--99.7 rule behind the
z-score, the IQR formula and box-plot interpretation, the pros and cons of mean, median and mode imputation, ML models that predict missing values (``XGB
imputer''), and leaky variables. This chapter completes those, with numbers.

\section{The normal curve and the 68--95--99.7 rule}
Many measurements (heights, test scores, measurement errors, averages of many things) follow a bell-shaped \textbf{normal distribution} with mean $\mu$ and
standard deviation $\sigma$ (Chapter 38). The \textbf{empirical rule}:
\begin{itemize}
\item about \textbf{68\%} of values lie within $\mu\pm1\sigma$,
\item about \textbf{95\%} within $\mu\pm2\sigma$ (exactly 95\% within $\pm1.96\sigma$),
\item about \textbf{99.7\%} within $\mu\pm3\sigma$.
\end{itemize}
\begin{center}
\begin{tikzpicture}
\begin{axis}[width=0.75\linewidth,height=4.2cm,axis lines=left,ytick=\empty,xtick={-3,-2,-1,0,1,2,3},
  xticklabels={$\mu{-}3\sigma$,$\mu{-}2\sigma$,$\mu{-}\sigma$,$\mu$,$\mu{+}\sigma$,$\mu{+}2\sigma$,$\mu{+}3\sigma$},ymin=0,ymax=0.45,
  tick label style={font=\scriptsize}]
\addplot[domain=-3.5:3.5,samples=80,thick,accent]{exp(-x^2/2)/sqrt(2*pi)};
\addplot[fill=accent,fill opacity=0.25,draw=none,domain=-1:1,samples=40]{exp(-x^2/2)/sqrt(2*pi)}\closedcycle;
\addplot[fill=accent,fill opacity=0.12,draw=none,domain=-2:2,samples=40]{exp(-x^2/2)/sqrt(2*pi)}\closedcycle;
\node[font=\scriptsize] at (axis cs:0,0.15) {68\%};
\node[font=\scriptsize] at (axis cs:1.5,0.06) {95\%};
\node[font=\scriptsize] at (axis cs:2.6,0.03) {99.7\%};
\end{axis}
\end{tikzpicture}
\end{center}
\begin{example}[: using the rule]
MBA percentage $\sim$ normal with mean 60, sd 5. About 68\% of students score 55--65; 95\% score 50--70; 99.7\% score 45--75. Of 1{,}000 students we expect
about 683, 954 and 997 in those ranges; only about 3 outside 45--75, so a score of 80 ($z = 4$) is genuinely rare. Also: by symmetry, 16\% score above 65
($\frac{100 - 68}{2}$) and 2.5\% above 70.
\end{example}

\section{z-scores}
\[ z = \frac{x - \mu}{\sigma} \]
says ``how many standard deviations from the mean''. It puts different variables on one scale and, for normal data, converts directly into rarity via the empirical rule.
\begin{example}[: z-score outlier check]
Daily applications to a job: mean 50, sd 8. A day with 76: $z = 26/8 = 3.25$: beyond 3, flagged (for a normal variable, $P(|Z|>3) = 0.27\%$).
\end{example}
\textbf{Limitations}: (1) assumes roughly normal data; for skewed data (salaries) many genuine values exceed 3; (2) \textbf{masking}: the outlier inflates the mean and sd it is
judged by, so in small samples an outlier may not even reach $|z| = 3$ (the 40 in Chapter 13 had $z = 2.4$).
\textbf{Modified z-score} (robust): $M_i = 0.6745\,\frac{x_i - \text{median}}{\text{MAD}}$, flag $|M_i| > 3.5$. The constant 0.6745 makes MAD comparable to $\sigma$ for normal data.
\begin{example}[: modified z]
Salaries 4, 5, 5, 6, 7, 9, 20: median 6, MAD 1. For 20: $M = 0.6745\times14/1 = 9.4$: clearly flagged, while the ordinary $z = (20 - 8)/5.54 = 2.2$ is not.
\end{example}

\section{IQR rule and box plots}
\[ \text{IQR} = Q_3 - Q_1,\qquad \text{lower fence} = Q_1 - 1.5\,\text{IQR},\qquad \text{upper fence} = Q_3 + 1.5\,\text{IQR}. \]
Points outside the fences are outliers (``far out'' beyond $3\times$IQR). It uses quartiles, so it is robust and does not assume normality; it is what box-plot
whiskers show (Chapter 31). For right-skewed data it flags many points on the high side only, which is correct behaviour (the high tail is long), but those
points are usually genuine.
\begin{example}[: IQR]
Data 3, 5, 7, 8, 9, 11, 13, 40 with $Q_1 = 6$, $Q_3 = 12$: IQR 6, fences $-3$ and 21: 40 is an outlier.
\end{example}

\section{Multivariate and model-based detection}
A point can look normal on every single feature but be unusual in combination (age 22 with 15 years of experience). Methods:
\textbf{Mahalanobis distance} $\sqrt{(\mathbf x - \boldsymbol\mu)^\top\Sigma^{-1}(\mathbf x - \boldsymbol\mu)}$ (a z-score that accounts for correlations; squared, it follows a
$\chi^2_d$ for normal data); \textbf{Isolation Forest} (Chapter 7); \textbf{Local Outlier Factor} (density relative to neighbours); \textbf{DBSCAN} noise points;
autoencoder reconstruction error.

\section{What to do with an outlier: a decision procedure}
\begin{enumerate}
\item \textbf{Is it impossible?} (negative age, percentage 105, price Rs.~1) $\to$ data error: correct from source or set to missing.
\item \textbf{Is it genuine but extreme?} (a 2-crore salary for a CTO) $\to$ keep; reduce its influence: log transform, winsorise, robust loss (Huber, MAE),
  robust scaler, tree models.
\item \textbf{Is it from a different population?} (bulk staffing agencies among recruiters) $\to$ segment or add an indicator.
\item \textbf{Is it the signal?} (fraud) $\to$ anomaly detection; never remove.
\end{enumerate}
Always decide using training data only and apply the same rule to test data; never drop test rows.

\section{Missing values: the full picture}
\subsection{Diagnosing the mechanism}
MCAR, MAR, MNAR (Chapter 13). Practical checks: compare the distributions of other variables for rows with and without the missing value (if they differ, it is not
MCAR); compare target rates for missing vs present; ask how the data were collected. Little's MCAR test exists but is rarely decisive.

\subsection{Dropping}
\begin{itemize}
\item \textbf{Listwise deletion} (drop rows): unbiased only under MCAR; loses data fast with many columns ($0.95^{12} = 54\%$ complete rows).
\item \textbf{Pairwise deletion} (use available pairs for each correlation): keeps more data but statistics can be inconsistent with each other.
\item \textbf{Dropping a column}: when mostly missing \emph{and} uninformative.
\item Never drop rows whose missingness is tied to the target (salary missing for every non-placed student).
\end{itemize}

\subsection{Mean, median and mode imputation: advantages and disadvantages}
\begin{center}\small
\begin{tabularx}{\linewidth}{l X X}
\toprule
Method & Advantages & Disadvantages\\
\midrule
Mean & simple; keeps the mean unchanged; fine for symmetric data, small missing fractions, MCAR & pulled by outliers; shrinks variance; weakens correlations; distorts
distribution (spike at the mean); ignores relationships with other features; biased under MAR/MNAR\\
Median & robust to outliers and skew (salary, price); simple & same variance shrinkage and correlation weakening; spike at the median\\
Mode & works for categoricals and discrete values & over-represents the most frequent category; can be arbitrary when several modes; hides the signal in
missingness (better: a separate ``Unknown'' category)\\
Constant ($-999$, 0) & preserves the information ``was missing'' for tree models & meaningless for linear models and distances unless paired with an indicator\\
\bottomrule
\end{tabularx}
\end{center}
\begin{example}[: variance shrinkage, simulated]
500 salaries, 20\% removed at random, true variance 29.2. After mean imputation the column's variance is 21.5: the filled values sit at the centre and pretend we are more
certain than we are.
\end{example}

\subsection{ML models that predict missing values}
The idea: treat the column with gaps as a \emph{target}, learn it from the other columns on rows where it is observed, then predict it where it is missing.
\begin{itemize}
\item \textbf{KNN imputer}: average of the $k$ most similar rows (scale first).
\item \textbf{Regression imputation}: linear model from the other columns (add noise for stochastic imputation).
\item \textbf{MICE / IterativeImputer}: cycle through every incomplete column, modelling each from all the others, for several rounds.
\item \textbf{MissForest}: IterativeImputer with random forests: handles non-linearity, interactions and mixed types.
\item \textbf{``XGB imputer''}: the same idea with gradient boosting (XGBoost, LightGBM, \code{HistGradientBoosting}), whose ability to accept NaN in the
  \emph{predictor} columns means rows with several gaps can still be used.
\end{itemize}
\begin{lstlisting}
from sklearn.experimental import enable_iterative_imputer   # noqa: enables the class
from sklearn.impute import IterativeImputer, KNNImputer
from sklearn.ensemble import RandomForestRegressor, HistGradientBoostingRegressor

knn  = KNNImputer(n_neighbors=5)
mice = IterativeImputer(random_state=0)                             # Bayesian ridge per column
miss_forest = IterativeImputer(estimator=RandomForestRegressor(n_estimators=50), max_iter=5)

def gbm_impute(df, col):                 # "XGB imputer" for one column
    obs = df[col].notna()
    feats = df.columns.drop(col)
    m = HistGradientBoostingRegressor().fit(df.loc[obs, feats], df.loc[obs, col])
    out = df[col].copy()
    out[~obs] = m.predict(df.loc[~obs, feats])
    return out
\end{lstlisting}
\begin{example}[: how much better? (simulated)]
Salary depends on experience; 20\% of salaries removed at random. RMSE of the imputed values: mean 5.81, median 5.79, KNN 1.69, MICE 1.64, MissForest 1.72,
gradient-boosting imputer 1.68. Model-based imputation is more than 3 times more accurate because it uses the relationship with experience.
\end{example}
Cautions: fit the imputer on training data only (inside a Pipeline); model-based imputations are ``too clean'' (no noise), which can understate uncertainty; for
inference prefer multiple imputation; add a missing indicator when missingness may be informative; for GBM models, native NaN handling is often as good and simpler.

\section{Leaky variables}
A \textbf{leaky variable} contains information about the target that would not be available at prediction time. Catalogue with detection tips:
\begin{center}\small
\begin{tabularx}{\linewidth}{l X X}
\toprule
Type & Example & How to catch it\\
\midrule
Outcome recorded after the event & salary for placement; ``offer letter generated'' for hiring; ``refund issued'' for fraud & ask ``when is this value created?'';
missingness perfectly tied to the target\\
Proxy of the target & ``number of interviews'' when predicting ``hired''; account status ``closed'' for churn & suspiciously high importance or AUC of one feature\\
Future information & features aggregated over the full period including after the prediction date & time-based validation drops sharply\\
Preprocessing on all data & scaler, imputer, TF-IDF, target encoding fitted before the split & always use Pipelines; fit transforms in folds\\
Duplicates / same entity & same candidate in train and test & group split by entity\\
ID or timestamp ordering & row ID correlated with label due to sorted export & drop IDs; check importance\\
\bottomrule
\end{tabularx}
\end{center}

\begin{practical}[: where this chapter shows up in ML]
z-scores standardise features (StandardScaler) and power anomaly alerts on metrics dashboards (``applications today are 4 sd below normal: check the pipeline'');
IQR rules drive winsorising; imputation choices change model accuracy directly; and most ``too good to be true'' offline results are leakage.
\end{practical}

\stuck{Kaggle Learn: Missing Values}{https://www.kaggle.com/code/alexisbcook/missing-values}
\section{Interview questions}

\begin{qa}{State the 68--95--99.7 rule and use it: MBA scores are normal with mean 60 and sd 5. What fraction score above 70? Between 55 and 65?}
\oneline{About 68\%, 95\% and 99.7\% of a normal variable lie within 1, 2 and 3 standard deviations of the mean; above 70 is $z > 2$, about 2.5\%; between 55 and 65 is
within $\pm1\sigma$, about 68\%.}
\why Above $\mu + 2\sigma$: $(100 - 95)/2 = 2.5\%$ by symmetry.
\end{qa}

\begin{qa}{How do you detect outliers with the z-score, and what are its weaknesses?}
\oneline{Standardise $z = (x - \mu)/\sigma$ and flag $|z| > 3$ (or 2.5); it assumes approximate normality and suffers from masking because the outlier inflates $\mu$ and
$\sigma$; robust alternatives are the modified z-score with median and MAD, or the IQR rule.}
\worked Salaries with 20: ordinary $z = 2.2$ (missed), modified $z = 9.4$ (caught).
\end{qa}

\begin{qa}{Write the IQR rule and explain why the multiplier is 1.5.}
\oneline{Outliers lie below $Q_1 - 1.5\,$IQR or above $Q_3 + 1.5\,$IQR; with 1.5 the fences fall at about $\pm2.7\sigma$ for normal data, flagging roughly 0.7\% of points,
a sensible threshold for ``unusual'' without assuming normality for the rule itself.}
\end{qa}

\begin{qa}{Would you remove outliers before training? Explain your decision process.}
\oneline{Only data errors; genuine extremes are information, so I reduce their influence (transforms, winsorising, robust losses, tree models) and treat
target-relevant anomalies as signal; decisions are made on training data and applied consistently.}
\end{qa}

\begin{qa}{Give the advantages and disadvantages of mean, median and mode imputation.}
\oneline{All are fast and simple; mean preserves the average but is outlier-sensitive; median is robust for skewed data; mode works for categories; all shrink variance, weaken
correlations, create artificial spikes, ignore other features and are biased unless data are MCAR.}
\end{qa}

\begin{qa}{\deep How can an ML model be used to impute missing values? Describe an ``XGB imputer''.}
\oneline{Treat the incomplete column as a target, train a model (KNN, linear, random forest, gradient boosting) on rows where it is observed using the other columns as features,
and predict it for rows where it is missing; with gradient boosting, the predictor columns may themselves contain NaNs, which the trees handle natively; iterating over all
incomplete columns gives MICE/MissForest-style imputation.}
\worked Simulated: RMSE 5.8 for mean imputation vs about 1.7 for model-based.
\pushes \emph{``Risks?''} Leakage if fitted on test data; using the target as an imputation feature for test rows (not available at prediction time); understated
uncertainty; extra complexity.
\end{qa}

\begin{qa}{\deep Should you use the target variable when imputing features?}
\oneline{For pure prediction, no: at prediction time the target is unknown, so an imputer that relied on it cannot be applied consistently and leaks; for statistical inference with
multiple imputation, including the outcome in the imputation model is actually recommended, because the analysis model will condition on it.}
\end{qa}

\begin{qa}{What is a leaky variable? Give three examples and how you would detect each.}
\oneline{A feature carrying information about the target that would not exist at prediction time; e.g. salary for placement (missingness perfectly tied to the label),
``number of interviews'' for hiring (a post-outcome proxy with extreme importance), and full-period aggregates for forecasting (validation collapses with a time split).}
\end{qa}

\begin{qa}{What is the Mahalanobis distance and why is it better than per-feature z-scores for multivariate outliers?}
\oneline{It measures distance from the mean in units that account for the covariance between features, so it flags combinations that are unusual given the correlations (age 22 with
15 years of experience) even when each value alone is normal.}
\end{qa}

\begin{qa}{Listwise vs pairwise deletion?}
\oneline{Listwise drops any row with a missing value (consistent but wasteful); pairwise uses all available data for each pair of variables (keeps more data but different statistics
use different subsets and a correlation matrix may not be valid).}
\end{qa}
